\documentclass[border=8pt,varwidth=175mm]{standalone}
% Task 024a -- shared figure style for the Matrix Operator Refactor benchmark figures.
%
% Every figure in this directory is a standalone LaTeX document that \input{}s this
% file and reads its data from tables/*.csv (generated by
% ../../scripts/figures_024a_prepare.jl).  Build with ../../scripts/figures_024a_build.sh.
%
% Palette: four categorical slots, validated colourblind-safe on a white surface
% (all-pairs mode: lightness band PASS, chroma floor PASS, worst CVD dE 8.4,
% worst normal-vision dE 16.3, all four >= 3:1 contrast).  Series are additionally
% separated by mark shape, so identity never rests on colour alone.
% Slots are assigned in fixed order and never cycled or recoloured by rank.

\usepackage[T1]{fontenc}
\usepackage[scaled]{helvet}
\renewcommand{\familydefault}{\sfdefault}
\usepackage{amsmath}
\usepackage{xcolor}
\usepackage{pgfplots}
\usepackage{pgfplotstable}
\usepgfplotslibrary{groupplots}
\usetikzlibrary{patterns}
\pgfplotsset{compat=1.18}

% ---- categorical slots (fixed order) -------------------------------------
\definecolor{fmSone}{HTML}{2A78D6}   % blue
\definecolor{fmStwo}{HTML}{EB6834}   % orange
\definecolor{fmSthree}{HTML}{199E70} % aqua
\definecolor{fmSfour}{HTML}{4A3AA7}  % violet
% ---- chrome and ink ------------------------------------------------------
\definecolor{fmInk}{HTML}{0B0B0B}
\definecolor{fmInkTwo}{HTML}{52514E}
\definecolor{fmMuted}{HTML}{898781}
\definecolor{fmGrid}{HTML}{E1E0D9}
\definecolor{fmAxis}{HTML}{C3C2B7}
% ---- reserved status colour: gates, rooflines, break-even rules ----------
\definecolor{fmCritical}{HTML}{D03B3B}

% Same reference-rule look as the \pgfplotsset fmrule style below, but usable from
% \draw inside an axis (needed on ybar axes, where an \addplot would become a bar).
\tikzset{fmruleline/.style={color=fmCritical, line width=0.7pt, dash pattern=on 3pt off 2pt}}

\newcommand{\fmtitle}[1]{{\bfseries\small\color{fmInk}#1}\par\vspace{2pt}}
\newcommand{\fmnote}[1]{\vspace{1pt}{\scriptsize\color{fmInkTwo}#1}\par}

\pgfplotsset{
  % recessive grid and axes, thin marks, no chart junk
  fmbase/.style={
    width=72mm, height=52mm,
    scale only axis,
    unbounded coords=jump,          % 'nan' in the tables is skipped, never drawn as 0
    every axis plot/.append style={line width=0.9pt, mark size=2.1pt},
    tick label style={font=\scriptsize, color=fmInkTwo},
    label style={font=\scriptsize, color=fmInkTwo},
    title style={font=\small\bfseries, color=fmInk, yshift=1pt},
    legend style={font=\scriptsize, draw=fmAxis, fill=white, fill opacity=0.92,
                  text opacity=1, inner sep=2pt, row sep=-1pt, /tikz/every even column/.append style={column sep=4pt}},
    legend cell align=left,
    grid=major,
    grid style={fmGrid, line width=0.35pt},
    axis line style={fmAxis, line width=0.5pt},
    tick style={fmAxis, line width=0.5pt},
    major tick length=2pt,
    minor tick length=1pt,
    log ticks with fixed point,
  },
  % series slots: colour + mark together (composite encoding)
  fmA/.style={color=fmSone,   mark=*},
  fmB/.style={color=fmStwo,   mark=square*},
  fmC/.style={color=fmSthree, mark=triangle*},
  fmD/.style={color=fmSfour,  mark=diamond*},
  % second channel of the same entity: same colour + mark, dashed, hollow
  fmalt/.style={dashed, mark options={fill=white}},
  % reference rules (break-even, roofline, memory gate) -- never a data series
  fmrule/.style={color=fmCritical, line width=0.7pt, dash pattern=on 3pt off 2pt,
                 mark=none, forget plot},
  % grouped bars: fmbaraxis on the axis (positions the groups, sets the width via
  % the tikz key), fmbars on each \addplot (flat fill, no outline, no marks).
  % 'bar width' lives in /tikz, so it must be spelled out at axis level.
  % log origin y=infty anchors log-scale bars to the axis floor; the pgfplots
  % default (10^0) would draw them from y=1 and misrepresent sub-1 values.
  fmbaraxis/.style={ybar, /tikz/bar width=3.4pt, enlarge x limits=0.14,
                    log origin y=infty},
  fmbars/.style={draw=none, mark=none},
}

\begin{document}

\fmtitle{Fig.\ 6 --- What the speedups cost: the $\chi$ channel, multi-thread BLAS, construction, and accuracy}

\begin{tikzpicture}
\begin{groupplot}[
  fmbase,
  width=66mm, height=44mm,
  group style={group size=2 by 2, horizontal sep=16mm, vertical sep=17mm},
]

% --- (a) Lamb-Helmholtz cost multiplier -----------------------------------
\nextgroupplot[title={(a) Lamb--Helmholtz cost multiplier}, xlabel={expansion order $P$},
               ylabel={LH step / $\varphi$-only step}, xmin=3, xmax=13, xtick={4,8,12},
               ymin=1, ymax=3.1, legend pos=north west, legend columns=2]
\addplot[fmA] table[x=P, y=concat_ratio, col sep=comma] {tables/fig06a_lh_ratio_cpu.csv};
\addlegendentry{concat, CPU}
\addplot[fmB] table[x=P, y=factored_ratio, col sep=comma] {tables/fig06a_lh_ratio_cpu.csv};
\addlegendentry{factored, CPU}
\addplot[fmC] table[x=P, y=precomputed_y_ratio, col sep=comma] {tables/fig06a_lh_ratio_cpu.csv};
\addlegendentry{precomp.-$y$, CPU}
\addplot[fmD] table[x=P, y=dense_ratio, col sep=comma] {tables/fig06a_lh_ratio_cpu.csv};
\addlegendentry{dense, CPU}
\addplot[fmA, fmalt] table[x=P, y=concat_ratio, col sep=comma] {tables/fig06a_lh_ratio_gpu.csv};
\addlegendentry{concat, H200}
\addplot[fmD, fmalt] table[x=P, y=dense_ratio, col sep=comma] {tables/fig06a_lh_ratio_gpu.csv};
\addlegendentry{dense, H200}

% --- (b) multi-thread BLAS degradation ------------------------------------
\nextgroupplot[title={(b) 64-thread BLAS vs single-thread BLAS}, xlabel={expansion order $P$},
               ylabel={BLAS-1 time / BLAS-64 time}, xmin=0.6, xmax=8.5,
               xtick={1,2,3,4,6,8}, ymode=log, ymin=0.2, ymax=3.5,
               ytick={0.25,0.5,1,2}, legend pos=south west, legend columns=1]
\addplot[fmrule] coordinates {(0.6,1) (8.5,1)};
\addplot[fmA] table[x=P, y=tinyparent_phi, col sep=comma] {tables/fig06b_blas_degradation.csv};
\addlegendentry{tiny, $\varphi$}
\addplot[fmA, fmalt] table[x=P, y=tinyparent_lh, col sep=comma] {tables/fig06b_blas_degradation.csv};
\addlegendentry{tiny, LH}
\addplot[fmB] table[x=P, y=smallconstp_phi, col sep=comma] {tables/fig06b_blas_degradation.csv};
\addlegendentry{small, $\varphi$}
\addplot[fmC] table[x=P, y=mediumconstp_phi, col sep=comma] {tables/fig06b_blas_degradation.csv};
\addlegendentry{medium, $\varphi$}
\node[anchor=north west, font=\tiny, color=fmCritical, align=left]
  at (rel axis cs:0.03,0.97) {below 1: the 64-thread\\BLAS is the slower one};

% --- (c) construction amortization ---------------------------------------
\nextgroupplot[title={(c) dense construction break-even}, xlabel={expansion order $P$},
               ylabel={recurring steps to amortize}, xmin=3, xmax=13, xtick={4,8,12},
               ymode=log, ymin=1, ymax=2e7,
               legend style={at={(0.5,-0.32)}, anchor=north, legend columns=2, draw=fmAxis}]
\addplot[fmA] table[x=P, y=cpu_min, col sep=comma] {tables/fig06c_break_even.csv};
\addlegendentry{CPU, best case}
\addplot[fmA, fmalt] table[x=P, y=cpu_max, col sep=comma] {tables/fig06c_break_even.csv};
\addlegendentry{CPU, worst case}
\addplot[fmB] table[x=P, y=gpu_min, col sep=comma] {tables/fig06c_break_even.csv};
\addlegendentry{H200, best case}
\addplot[fmB, fmalt] table[x=P, y=gpu_max, col sep=comma] {tables/fig06c_break_even.csv};
\addlegendentry{H200, worst case}

% --- (d) accuracy vs cost ------------------------------------------------
\nextgroupplot[title={(d) accuracy vs cost (H200)}, xlabel={complete recurring step (ms)},
               ylabel={max direct potential error}, xmode=log, ymode=log,
               every axis plot/.append style={only marks, mark size=1.9pt},
               ymin=3e-12, ymax=3e-5,
               ytick={1e-11,1e-10,1e-9,1e-8,1e-7,1e-6,1e-5},
               yticklabels={$10^{-11}$,$10^{-10}$,$10^{-9}$,$10^{-8}$,$10^{-7}$,$10^{-6}$,$10^{-5}$},
               legend style={at={(0.5,-0.32)}, anchor=north, legend columns=3, draw=fmAxis}]
\addplot[fmA] table[x=step_ms, y=error, col sep=comma] {tables/fig06d_potential_float64_p4.csv};
\addlegendentry{F64 $P{=}4$}
\addplot[fmB] table[x=step_ms, y=error, col sep=comma] {tables/fig06d_potential_float64_p8.csv};
\addlegendentry{F64 $P{=}8$}
\addplot[fmC] table[x=step_ms, y=error, col sep=comma] {tables/fig06d_potential_float64_p12.csv};
\addlegendentry{F64 $P{=}12$}
\addplot[fmA, fmalt] table[x=step_ms, y=error, col sep=comma] {tables/fig06d_potential_float32_p4.csv};
\addlegendentry{F32 $P{=}4$}
\addplot[fmB, fmalt] table[x=step_ms, y=error, col sep=comma] {tables/fig06d_potential_float32_p8.csv};
\addlegendentry{F32 $P{=}8$}
\addplot[fmC, fmalt] table[x=step_ms, y=error, col sep=comma] {tables/fig06d_potential_float32_p12.csv};
\addlegendentry{F32 $P{=}12$}

\end{groupplot}
\end{tikzpicture}

\vspace{5mm}

\fmnote{\textbf{Machines / regimes:} (a,c,d) task 024 --- CPU \texttt{m12-2-17} EPYC 7763 (BLAS 1 and 64,
\texttt{julia\_threads=1}), GPU \texttt{m13h-1-1} H200. (b) task 019b, \texttt{m12-1-7} EPYC 7763, whole-slab
concat host stage, Float64.
\textbf{Data:} \texttt{operator\_ab\_benchmark/summary/\{case\_rankings,construction\_amortization\}.csv},
\texttt{operator\_ab\_benchmark/raw/cuda\_*.csv},
\texttt{smallp\_fallback\_layout/m12-1-7/crossover\_stage\_blas\{1,64\}.csv}
(tables \texttt{fig06a\_*} \dots \texttt{fig06d\_*}).
\textbf{Reads as:} carrying the Lamb--Helmholtz $\chi$ channel costs 1.3$\times$ at $P=4$ rising to
2.1--2.8$\times$ at $P=12$, and it is the axis that most often flips the strategy choice in Fig.\ 5. Handing the
operators a 64-thread BLAS \emph{loses} up to 3--4$\times$ at small route counts, which is why the standing rule
is to pin BLAS to one thread inside the operators rather than to add a fallback path. The dense strategy's
construction amortizes in 4--220 steps across the measured large-$n$ CPU cases
(4--35 for uniform Float64, 6--109 for clustered Float64, and 52--220 for uniform Float32), but needs
$5\times10^2$--$7\times10^5$ steps on the H200 --- the reason precomputed-$y$ is the recommended GPU default
despite dense winning many individual steps.
\textbf{Scope:} (a,c,d) use the complete recurring \texttt{fmm!} step, far field \emph{and} nearfield (the
latter inside L2B); panel (b) is the far-field M2L host stage alone. Because the constant-$P$ stencil promotes
more pairs to the far field as $P$ grows, the nearfield share of (a,c,d) shrinks with $P$ --- at $n=20000$ the
direct-pair count falls from 227\,632 at $P=4$ to 39\,496 at $P=12$ --- so the two costs trade against each
other along the $P$ axis.
\textbf{Caveats:} \emph{panel (d) is not an accuracy-vs-tolerance study.} Every 024 case ran a single fixed
\texttt{ConstantPAnalyticStencil} tolerance, so the only accuracy axis available is Float32 vs Float64 at
each $P$; the spread along $x$ within one series is the four strategies and the three problem sizes, not an
accuracy knob. No stencil-tolerance sweep exists anywhere in the project's data --- tolerance-vs-cost remains
unmeasured and is carried to the 019a review as a limitation. In (c), $P=12$ has no GPU point because dense is
not a steady-state winner against precomputed-$y$ there; the CPU worst case at each $P$ is dominated by the
$n=150$ cases (up to 18\,930 steps). Panel (b) gaps at $P\in\{1,3,6,8\}$ for the medium configuration are
unrun points.}

\end{document}
